\subsection{Reconstruction with fixed optics}
\label{sec:decoders}

\begin{figure}[t]
\centering
\includegraphics[width=\linewidth]{reconstruction.pdf}
\caption{\textbf{Reconstruction is evaluated after the plate is fixed.} TV and PnP use the alternating direction method of multipliers (ADMM) with different prior updates. The feed-forward route uses a regularized solve followed by a 3D U-Net whose weights are trained separately for each plate. Full updates are in the supplement.}
\label{fig:reconstruction}
\end{figure}

We evaluate three iterative methods and a feed-forward decoder using the operator $\bA_\phi$ of each fixed plate (\cref{fig:reconstruction}). TV combines a quadratic measurement term with isotropic 3D total variation and $[0,1]$ intensity bounds \cite{rudin1992}. PnP uses denoisers as implicit priors within ADMM \cite{venkatakrishnan2013,kamilov2023pnp}. Following multi-slice fusion \cite{majee2021}, pretrained PnP applies a natural-image DnCNN independently along three orthogonal slice orientations of a shared volume \cite{zhang2017}. For slice denoising, neuron-trained PnP uses a 2D network trained on neuronal volumes.

The TV regularizer's form and PnP denoiser weights remain fixed across plates. For these iterative methods, scalar parameters are selected separately for each plate and acquisition condition using held-out validation volumes. Feed-forward decoders share an architecture, with weights fine-tuned separately for each fixed plate. Full solver updates, denoiser configurations and training protocols are provided in the supplement.
